\documentclass[11pt,a4paper]{article}
\usepackage{fontspec}
\usepackage{geometry}
\usepackage{amsmath,amssymb,mathtools}
\usepackage{booktabs}
\usepackage{enumitem}
\usepackage{xcolor}
\usepackage{titlesec}
\usepackage{fancyhdr}
\usepackage[unicode=true]{hyperref}
\usepackage[most]{tcolorbox}
\usepackage{lastpage}
\usepackage{microtype}
\usepackage{xeCJK}
\geometry{top=22mm,bottom=24mm,left=24mm,right=24mm,headheight=14pt}
\definecolor{ink}{HTML}{172033}\definecolor{accent}{HTML}{315A8A}\definecolor{accentlight}{HTML}{EAF1F8}\definecolor{rulegray}{HTML}{D7DEE8}\definecolor{muted}{HTML}{5C6678}
\setmainfont{DejaVu Serif}\setsansfont{DejaVu Sans}\setmonofont{DejaVu Sans Mono}\setCJKmainfont{Noto Serif CJK SC}\setCJKsansfont{Noto Sans CJK SC}
\hypersetup{colorlinks=true,linkcolor=accent,urlcolor=accent,citecolor=accent,pdftitle={模 3 数列的半周期对称性},pdfsubject={数学形式化与算法审计},pdfkeywords={递推, 模3, 反周期性, 傅里叶, 伴随矩阵}}
\pagestyle{fancy}\fancyhf{}\fancyhead[L]{\small\color{muted}模 3 半周期对称性}\fancyhead[R]{\small\color{muted}形式化与审计}\fancyfoot[C]{\small\color{muted}\thepage\,/\,\pageref*{LastPage}}\renewcommand{\headrulewidth}{0.3pt}\renewcommand{\footrulewidth}{0pt}
\titleformat{\section}{\Large\bfseries\color{ink}}{\thesection.}{0.6em}{}[\vspace{0.2em}\color{rulegray}\titlerule]\titleformat{\subsection}{\large\bfseries\color{accent}}{\thesubsection}{0.6em}{}\titlespacing*{\section}{0pt}{2.2ex plus .5ex minus .2ex}{1.3ex}\titlespacing*{\subsection}{0pt}{1.7ex plus .4ex minus .2ex}{0.8ex}
\setlist[itemize]{leftmargin=1.6em,itemsep=0.25em,topsep=0.4em}\setlist[enumerate]{leftmargin=1.8em,itemsep=0.25em,topsep=0.4em}\setlength{\parindent}{0pt}\setlength{\parskip}{0.55em}\allowdisplaybreaks
\newtcolorbox{summarybox}{colback=accentlight,colframe=accent,boxrule=0.6pt,arc=2mm,left=4mm,right=4mm,top=3mm,bottom=3mm,before skip=1em,after skip=1em}\newcommand{\F}{\mathbb F}\newcommand{\Z}{\mathbb Z}\newcommand{\concat}{\mathbin{\Vert}}
\begin{document}
\begin{titlepage}\thispagestyle{empty}\vspace*{22mm}{\color{accent}\rule{\textwidth}{1.2pt}}\\[15mm]{\fontsize{27}{33}\selectfont\bfseries\color{ink}模 3 数列的半周期对称性\par}\vspace{9mm}{\Large\color{accent}数学形式化与算法审计\par}\vspace{13mm}\begin{summarybox}本文严格区分三个常被混淆的现象：循环取负、半周期反周期性以及反转。文中证明非零本原词的唯一相反移位，给出组合、频谱和矩阵刻画，并把这些结论与模递推的穷尽分析器联系起来。\end{summarybox}\vfill{\large\color{muted}仓库: \texttt{Modular-recurrence-symetry-analyzer-}\par}{\large\color{muted}汇编 PDF 版本 - 2026 年 8 月 5 日\par}\vspace{12mm}{\color{accent}\rule{\textwidth}{1.2pt}}\end{titlepage}
\renewcommand{\contentsname}{目录}\tableofcontents\clearpage
\section{目的}
本文对仓库 \href{https://github.com/59200/k-bonacci-pisano-finder}{\texttt{59200/k-bonacci-pisano-finder}} 中以“对 3 互补的周期”和“切成两半后自互补”等非正式说法描述的现象进行形式化。

不变量框架是有限域 \(\F_3=\{0,1,2\}\) 上的周期词，其加法对合为 \(\nu(x)=-x\pmod3\)，即 \(\nu(0)=0\)、\(\nu(1)=2\)、\(\nu(2)=1\)。因此，“对 3 的补数”应改称 \textbf{模 3 加法逆元} 或 \textbf{模 3 取负补元}。所选整数代表的十进制和不总是 3，因为 0 被映到 0。
\section{三个不同的性质}
设 \(w=(w_0,\ldots,w_{T-1})\in\F_3^T\) 是本原周期为 \(T\) 的词。
\subsection{相反周期对}
若对所有 \(n\) 都有 \(v_n=-w_n\pmod3\)，则周期 \(w\) 与 \(v\) 构成相反对。它们可以属于两个不同的循环。
\subsection{半周期反周期性}
若 \(T=2h\) 且对所有 \(n\) 有 \(w_{n+h}=-w_n\pmod3\)，则称 \(w\) 具有\textbf{半周期反周期性}。此时 \(w=H\concat(-H)\)，其中 \(H=(w_0,\ldots,w_{h-1})\)。例如 \(01120221=0112\concat0221\)、\(011022=011\concat022\)、\(01220211=0122\concat0211\)。
\subsection{后半段反转}
运算 \(\mathcal R(a_0,\ldots,a_{h-1})=(a_{h-1},\ldots,a_0)\) 与取负相互独立。在斐波那契示例中，\(\mathcal R(0221)=1220\)。因此，\(1220\) 是 \(0112\) 的\textbf{反转对跖词}，而不是原来的后半段。
\section{唯一相反移位定理}
\begin{summarybox}\textbf{定理.} 设 \(w\) 是周期为 \(T\) 的非零本原词。若存在移位 \(r\)，使得对所有 \(n\) 都有 \(w_{n+r}=-w_n\)，则 \(T\) 为偶数，且 \(r\equiv T/2\pmod T\)。\end{summarybox}
\subsection*{证明}
连续应用两次移位得到 \(w_{n+2r}=w_n\)。本原性推出 \(T\mid2r\)。由于非零的 \(\F_3\) 词不可能满足 \(w=-w\)，故 \(r\not\equiv0\pmod T\)。循环群 \(\Z/T\Z\) 中唯一的非平凡二阶元仅在 \(T\) 为偶数时存在，并等于 \(T/2\)。\hfill\(\square\)
\subsection*{推论}
因此，非零本原周期只有两种情形：\begin{enumerate}\item 取负后落到另一个循环；\item 在取负下不变，此时必然具有半周期反周期性。\end{enumerate} 这一区分澄清了 \(00111201,00222102\) 这类相反循环对与 \(011022\) 这类内在反周期之间的差别。
\section{组合推论}
若 \(w=H\concat(-H)\)，则 1 与 2 的出现次数相等，0 的出现次数为偶数，且 \(\sum_{n=0}^{T-1}w_n=0\)（在 \(\F_3\) 中）。这些条件必要但不充分。取中心提升 \(\widetilde0=0\)、\(\widetilde1=1\)、\(\widetilde2=-1\)，生成多项式分解为 \[W(X)=\sum_{n=0}^{2h-1}\widetilde w_nX^n=(1-X^h)\sum_{n=0}^{h-1}\widetilde w_nX^n.\]
\section{傅里叶刻画}
令 \(\widehat w(k)=\sum_{n=0}^{2h-1}\widetilde w_ne^{-2\pi i kn/(2h)}\)。则 \(w_{n+h}=-w_n\) 当且仅当所有偶数指标 \(k\) 都满足 \(\widehat w(k)=0\)。正向证明中出现因子 \(1-(-1)^k\)，它在 \(k\) 为偶数时为零。反之，若所有偶模均消失，则该词属于奇模张成的子空间，移位 \(h\) 在此子空间上作用为 \(-I\)。这是精确的频谱判据，而非两半的目测比较。
\section{模 3 线性数列}
考虑 \(k\) 阶递推 \[u_{n+k}=c_0u_n+\cdots+c_{k-1}u_{n+k-1}\pmod3.\] 状态为 \(s_n=(u_n,\ldots,u_{n+k-1})^{\mathsf T}\)，演化为 \(s_{n+1}=Ms_n\)，其中 \(M\) 是伴随矩阵。
\subsection{纯周期性}
\(M\) 的行列式在符号意义下等于 \(c_0\)。在 \(\F_3\) 上，\(M\) 可逆当且仅当 \(c_0\neq0\)。此时每条状态轨道都是纯周期的；若 \(c_0=0\)，则可能出现必须与循环区分的前周期。
\subsection{循环的取负}
线性性给出 \(M(-s)=-M(s)\)，所以取负把每个循环映到长度相同的循环。对于非零循环只有两种情形：两个不同循环被取负互换，或同一循环在取负下不变并必然半周期反周期。零循环是唯一的退化例外：它逐点固定，不给出非平凡的本原反周期。
\subsection{单个种子的判据}
对种子 \(s_0\)，关系 \(M^hs_0=-s_0\) 蕴含任意线性输出都满足 \(u_{n+h}=-u_n\)。当 \(s_0\neq0\) 且循环本原时，这就是上述半周期反周期。
\subsection{一致判据}
所有种子都满足同一反周期关系 \(h\) 当且仅当 \(M^h=-I\)。等价地，最小多项式满足 \(\mu_M(X)\mid X^h+1\)（在 \(\F_3[X]\) 中），从而 \(M^{2h}=I\)。
\section{斐波那契示例}
对 \(M=\begin{pmatrix}0&1\\1&1\end{pmatrix}\)（在 \(\F_3\) 上），可得 \(M^4=-I\)、\(M^8=I\)。种子 \((0,1)\) 生成周期 \(01120221=0112\concat(-0112)\)。其后半段反转为 \(\mathcal R(0221)=1220\)。十进制读取下的巧合 \(1220=L_{14}+F_{14}=2F_{15}\) 只是附加编码恒等式，并非反周期性的一般推论。
\section{推广到任意模 3 数列}
分析给定周期词不需要递推假设。对任一本原周期，都可以计算其旋转循环、加法逆元、仿射对称 \(w_{n+r}=aw_n+b\)、反转对称 \(w_{r-n}=aw_n+b\)、可能的反周期性以及奇偶傅里叶模。只有在给定线性递推时才需要矩阵理论。
\section{推广到一般模数}
在 \(\Z/m\Z\) 上，自然对合仍是 \(x\mapsto-x\)，定义 \(w_{n+h}=-w_n\pmod m\) 仍然有效。当 \(m=2\) 时，取负等于恒等映射，反周期性退化为普通周期性；当 \(m>2\) 时，二者仍有实质区别。“对 \(m\) 的补数”依赖整数代表的选择，不应作为代数定义；不变量术语是“模 \(m\) 加法逆元”。
\section{算法审计}
\(k\) 阶模 \(m\) 递推的周期必须在有限状态空间 \((\Z/m\Z)^k\) 上检测，而不能在任意长前缀中寻找重复块。脚本使用：基 \(m\) 的紧凑状态编码；对单个种子采用 Brent 算法，时间 \(O(\mu+\lambda)\)、空间 \(O(1)\)；对所有种子进行精确函数图分解，时间 \(O(m^k)\)；在系数族之间并行；分别分类相反循环与反周期循环；并验证 \(M^h=-I\)。无需 Mathematica 内核或外部封装。
\section{可复现的扫描结果}
对仓库中十五个命名族的扫描，在 tritetranacci 族中恢复了长度为 24 和 8 的不同相反循环对，以及反周期 \(011022\)、\(01220211\) 和 \(12\)。第二次扫描覆盖 2 至 6 阶的 \(119\) 个非零二元系数向量：\(\sum_{k=2}^{6}(2^k-1)=119\)。结果为：13 个族满足全局判据 \(M^h=-I\)，88 个族至少含一个非平凡半周期反周期循环，96 个族至少含一对不同的相反循环。结果可由仓库中的 JSON 和 CSV 文件复现。
\section*{结论}\addcontentsline{toc}{section}{结论}
模 3 取负在周期词和线性递推循环上构成结构性对合。对非零本原词而言，在该对合下保持不变只能通过唯一的半周期移位发生。同一性质也可由生成多项式的分解、偶数傅里叶模的消失，或在线性框架中由矩阵关系 \(M^h=-I\) 等价识别。因此，算法分析能够清楚地区分循环的内在对称与不同相反循环之间的简单配对。
\end{document}